\section{Proof of the Cycle-Rank Upper Bound}\label{app:low-functional}

We restate the result to be proved.

\medskip
\noindent\textbf{Lemma~\ref{lem:low-functional} (Cycle rank gives an upper
bound for long paths).}
Let \(\keepp<1\). For every \(\thr>0\) there is an integer
\(\plen=\plen(\degb,\keepp,\thr)\) such that every finite graph \(H\) with
maximum degree at most \(\degb\) satisfies
\[
  \Lpath{\plen}{H}{\keepp}
  \le
  \thr+2\ballsz{\plen}\,\frac{\cyc{H}}{|V(H)|} .
\]

\medskip
The proof first handles forests and then reduces a general graph to a forest
by removing \(\cyc{H}\) edges. Besides the symbols of
Sections~\ref{subsec:notation} and~\ref{sec:difficult-graphs}, this appendix
uses two new symbols, the forest extinction bound \(\extp{\plen}\) and the
separator set \(\sep\).

\subsection{From the forest case to a general graph}

\begin{lemma}[Uniform extinction in finite forests]\label{lem:forest-extinction}
Let \(\keepp<1\). There is a function \(\extp{\plen}\) of the path length,
depending only on \(\degb\) and \(\keepp\), with
\[
  \extp{\plen}\longrightarrow0
  \qquad(\plen\to\infty),
\]
such that every finite forest \(H\) with maximum degree at most \(\degb\)
satisfies
\begin{equation}\label{eq:forest-extinction}
  \Lpath{\plen}{H}{\keepp}\le \extp{\plen}.
\end{equation}
\end{lemma}

\begin{proof}[Proof of Lemma~\ref{lem:low-functional}]
Choose a spanning tree in every connected component of \(H\). Their union is a
spanning forest \(H'\subseteq H\), and we call the other edges of \(H\) extra
edges. By the definition of cycle rank, there are exactly
\[
  |E(H)\setminus E(H')|=\cyc{H}
\]
extra edges.

Choose a vertex uniformly at random and keep each edge of \(H\) with
probability \(\keepp\). If the kept graph contains a path of length \(\plen\)
from the chosen vertex, then one of two cases holds:
\begin{enumerate}[label=(\arabic*)]
\item the path uses only edges of \(H'\);
\item the path uses at least one extra edge.
\end{enumerate}

In case~(1), the path lies in \(\perc{H'}{\keepp}\). Since the forest \(H'\)
has the same vertices as \(H\), the chosen vertex is also uniform in \(H'\), so
by Lemma~\ref{lem:forest-extinction}, case~(1) has probability at most
\(\extp{\plen}\).

In case~(2), we only count which vertices can be the start and ignore whether
the edges are kept. If a path of length \(\plen\) from a vertex \(v\) uses an
extra edge \(e\), then \(v\) is at distance at most \(\plen\) from one endpoint
of \(e\). So \(v\) lies in one of the two balls of radius \(\plen\) around the
endpoints of \(e\), which contain at most \(2\ballsz{\plen}\) vertices. Since
there are \(\cyc{H}\) extra edges, at most \(2\ballsz{\plen}\cyc{H}\) vertices
are possible starts. The start is uniform among the \(|V(H)|\) vertices, so
case~(2) has probability at most
\[
  2\ballsz{\plen}\,\frac{\cyc{H}}{|V(H)|}.
\]
Adding the two cases gives
\begin{equation}\label{eq:low-functional-explicit}
  \Lpath{\plen}{H}{\keepp}
  \le
  \extp{\plen}
  +2\ballsz{\plen}\,\frac{\cyc{H}}{|V(H)|}.
\end{equation}
Given \(\thr>0\), choose \(\plen\) so large that \(\extp{\plen}\le\thr\); this
is possible because \(\extp{\plen}\to0\), and \(\plen\) depends only on
\(\degb,\keepp,\thr\). Then~\eqref{eq:low-functional-explicit}
gives~\eqref{eq:low-functional}.
\end{proof}

Only the property \(\extp{\plen}\to0\) is used here, and the exact formula for
\(\extp{\plen}\) does not matter. It remains to prove the forest estimate.

\subsection{Uniform extinction in finite forests}

\begin{proof}[Proof of Lemma~\ref{lem:forest-extinction}]
\emph{Cutting the forest.} Fix an integer \(N\ge2\). Root each tree of \(H\) at
any vertex, and go through the vertices from the leaves toward the root, so
that every vertex comes after all vertices below it. At each vertex, look at
the vertices of its subtree, itself included, that are not yet assigned. If
there are at least \(N\) of them, put the vertex into a set \(\sep\) and
assign all of them to it. Since each vertex of \(\sep\) gets at least \(N\)
vertices and different vertices of \(\sep\) get disjoint parts, we obtain
\begin{equation}\label{eq:separator-size}
  |\sep|\le\frac{|V(H)|}{N}.
\end{equation}
After this process, every connected component of \(H-\sep\) has fewer than
\(N\) vertices.

\emph{Small trees.} Define
\[
  g(N)
  :=\sup_{(C,u)}\sum_{v\in V(C)}\keepp^{\dist_C(u,v)},
\]
where the supremum is over all trees \(C\) with fewer than \(N\) vertices and
maximum degree at most \(\degb\), and all vertices \(u\) of \(C\). Since at most
\(\min\{N,\degb(\degb-1)^{t-1}\}\) vertices are at distance \(t\ge1\) from
\(u\), we have
\begin{equation}\label{eq:g-bound}
  g(N)
  \le
  1+\sum_{t\ge1}\min\{N,\degb(\degb-1)^{t-1}\}\keepp^t ,
  \qquad
  \frac{g(N)}{N}\longrightarrow0\quad(N\to\infty).
\end{equation}
To see the limit, note that if \(\degb=2\), the sum is bounded by a constant.
If \(\degb\ge3\), we split the sum at \(t=\lceil\log_{\degb-1}N\rceil\); then
the whole sum is at most a constant times
\(\log N+N^{1+\log_{\degb-1}\keepp}\), and \(1+\log_{\degb-1}\keepp<1\) because
\(\keepp<1\).

\emph{Long paths reach \(\sep\).} Suppose \(\perc{H}{\keepp}\) contains a path
of length \(N\) from a vertex \(v\). This path has \(N+1\) vertices, so it does
not fit into one component of \(H-\sep\), and hence \(v\) is joined to \(\sep\)
by kept edges. A vertex \(v\notin\sep\) can reach \(\sep\) only through an
edge \(sa\) with \(s\in\sep\) and \(a\notin\sep\), where \(v\) lies in the
component of \(a\) in \(H-\sep\). Since paths in a forest are unique, this
happens with probability at most \(\keepp^{\dist(a,v)}\). There are at most
\(\degb|\sep|\) such edges \(sa\), and for each of them the sum over \(v\) is
at most \(g(N)\), so after adding the vertices of \(\sep\) themselves we get
\[
  \sum_{v\in V(H)}
  \Prx{v\text{ is joined to }\sep\text{ in }\perc{H}{\keepp}}
  \le |\sep|+\degb|\sep|g(N).
\]
Dividing by \(|V(H)|\), using~\eqref{eq:separator-size}, and taking
\(N=\plen\) for \(\plen\ge2\), we obtain~\eqref{eq:forest-extinction} with
\[
  \extp{\plen}
  :=\frac{1+\degb\,g(\plen)}{\plen}
  \quad(\plen\ge2),
  \qquad
  \extp{1}:=1,
\]
and \(\extp{\plen}\to0\) by~\eqref{eq:g-bound}; for \(\plen=1\) the bound is
trivial.
\end{proof}
